% !TeX root = ../教材题库_全部条目_转写工作稿.tex
% Source: XJ-B1-C04-SU-P150-Y1226
% 第4章 / 小结与复习 小结与复习 / 教材 p.150 / PDF 第157页
% OCR未校对：公式与文字均可能有误；未完成数学TeX转写；禁止直接发给学生
% 原题及相邻原解截图（核对用）：../题目原图/XJ-B1-C04-SU-P150-Y1226-01.webp
% 本文件只在转写完成后作为题干使用；OCR 条目尚不是可靠题干。
\begin{bankitem}{第4章 小结与复习 第5题}{第4章 / 小结与复习 小结与复习 / 教材 p.150 / PDF 第157页}{XJ-B1-C04-SU-P150-Y1226}{OCR未校对：公式与文字均可能有误；未完成数学TeX转写；禁止直接发给学生}
\begin{ocrdraft}
[教材p.150 / PDF第157页]
\par
5．用函数的观点回顾方程，解方程可以化归为函数零点的计算．对于大量
\par
没有求解公式的方程，化为求函数零点后，可以用二分法或其他近似方法求函
\par
数零点的近似值．用二分法求函数零点的条件是什么？操作步骤如何进行？
\end{ocrdraft}
\end{bankitem}
